\makeatletter
\def\input@path{{./}{../}{../../}{../../../}{../../../../}{preamble/}{../preamble/}{../../preamble/}{../../../preamble/}{../../../../book2026-2027/preamble/}}
\makeatother

\input{preamble_exams.tex}
\usepackage{titling}
\setlength{\droptitle}{-2em}
\pretitle{\par\begingroup\centering\LARGE}
\posttitle{\par\endgroup}
\preauthor{}\postauthor{}
\predate{}\postdate{}
\maketitlehookd{\vspace{-1em}}
\linespread{1.12}
\setlength{\parskip}{5pt}
\newgeometry{margin=0.65in,top=0.55in,bottom=0.65in}
\raggedbottom

\title{Grade 11 Term 2 Home Practice -- C2\\\large Point estimates and confidence levels\\[0.25em]
    \normalsize Global Economics -- 2026--2027}
\author{}
\date{}

\newcommand{\EstimationTask}{Identify the point estimate and each confidence interval. State the population parameter being estimated, and explain what the point estimate and the interval estimate or estimates communicate about it.}

\newcommand{\EstimationTable}[1]{%
    \noindent
    \begin{tabularx}{\linewidth}{@{}>{\raggedright\arraybackslash}p{0.17\linewidth}*{2}{>{\raggedright\arraybackslash}X}@{}}
        \toprule
        \textbf{Quantity} & \textbf{Calculation and result} & \textbf{Meaning in context} \\
        \midrule
        #1
        \bottomrule
    \end{tabularx}\par
    \vspace{4mm}
}

\newlist{answeroptions}{enumerate}{1}
\setlist[answeroptions]{label=\Alph*., leftmargin=4em, topsep=0pt, itemsep=2pt, parsep=0pt}

\begin{document}
\maketitle

\setcounter{tocdepth}{1}
\setcounter{secnumdepth}{2}
\phantomsection
\label{toc}
\tableofcontents

\section*{Evaluated criteria}
\addcontentsline{toc}{section}{Evaluated criteria}
\begin{description}
    \item[\textbf{C2}] Compares point and interval estimation and interprets confidence levels through their long-run coverage meaning, distinguishing this from the probability that a particular calculated interval contains the parameter.
\end{description}

\textbf{How to practise.} Try each task before opening its blue solution link. The arrows return to the contents, the section's problem list, or the problem. In Task 1, interpret the estimates reported from the study samples. In Task 2, use the evidence to choose one answer.

\section{Introduction}
\label{sec:introduction}

A survey estimates that cyclists spend an average of 40 minutes on a weekly trip to the market. The sample mean of 40 minutes is a \emph{point estimate}. The 95\% confidence level sets the critical value at $z^*=1.960$. With a standard error of $4/1.960$ minutes, the margin of error (MOE) is $1.960(4/1.960)=4$ minutes. The interval $[40-4,40+4]=[36,44]$ is 8 minutes wide and gives a range of plausible values for the population mean under the sampling model. A fresh sample could give a different sample mean and different endpoints, while the population mean stays fixed.

\clearpage

\section{Notation Reference and Practice Procedure}
\label{sec:notation-and-procedure}

For the normal procedures used here, let $n$ be the sample size and $\sigma$ be the known population standard deviation. MOE stands for \emph{margin of error}, the distance from the sample mean to either endpoint.

\begin{center}
\begin{tabular}{@{}ll@{}}
    \toprule
    Notation & Meaning \\
    \midrule
    $\mu$ & Population mean, the fixed parameter of interest \\
    $\bar{x}$ & Sample mean, a point estimate of $\mu$ \\
    $\mathrm{SE}=\sigma/\sqrt n$ & Standard error of $\bar{x}$ for independent observations \\
    $z^*$ & Standard normal critical value for the confidence level \\
    $\mathrm{MOE}=z^*\mathrm{SE}$ & Margin of error at the stated confidence level \\
    $[\bar{x}-\mathrm{MOE},\bar{x}+\mathrm{MOE}]$ & Two-sided confidence interval, with width $2\mathrm{MOE}$ \\
    \bottomrule
\end{tabular}
\end{center}

Use the same units for $\mu$, $\bar{x}$, $\mathrm{SE}$, $\mathrm{MOE}$, the endpoints, and the width.

\subsection{Choosing the critical value}

Write the confidence level $C$ as a decimal. On the standard normal curve, the central area from $-z^*$ to $z^*$ is $C$, leaving $(1-C)/2$ in each tail. The area to the left of $z^*$ is then $C+(1-C)/2=(1+C)/2$, so
\[
    z^*=\Phi^{-1}\!\left(\frac{1+C}{2}\right).
\]
\noindent
\begin{minipage}[t]{0.56\linewidth}
    \vspace{0pt}
    $\Phi$ gives the area to the left of a value on the standard normal curve. Its inverse, $\Phi^{-1}$, takes an area and finds the value with that much area to its left. For 95\% confidence, each tail has area $0.025$, so look for $0.975$ to the left of $z^*$.

    \smallskip
    Choose $z^*$ from the confidence level and calculate $\bar{x}$ from the sample data. Higher confidence requires a larger $z^*$ to enclose more area. With the same standard error, this gives a larger $\mathrm{MOE}$ and a wider interval.
\end{minipage}
\hfill
\begin{minipage}[t]{0.40\linewidth}
    \vspace{0pt}
    \centering
    \begin{tabular}{@{}lc@{}}
        \toprule
        Confidence level & \shortstack{Standard normal\\critical value $z^*$} \\
        \midrule
        80\% & 1.282 \\
        90\% & 1.645 \\
        95\% & 1.960 \\
        97.5\% & 2.241 \\
        99\% & 2.576 \\
        99.9\% & 3.291 \\
        \bottomrule
    \end{tabular}
    \par\smallskip
    \small Two-sided values rounded to three decimals.
\end{minipage}

\smallskip
\noindent
\begin{minipage}{\linewidth}
    \centering
    \begin{tikzpicture}
    \begin{axis}[
        width=0.82\linewidth,height=4.8cm,axis x line=bottom,axis y line=none,
        xmin=-4.2,xmax=4.2,ymin=0,ymax=0.43,
        xlabel={$z$},xlabel style={at={(axis description cs:1,0)},anchor=west,font=\small},
        xtick={-1.96,0,1.96},xticklabels={$-z^*=-1.960$,$0$,$z^*=1.960$},
        tick label style={font=\small},samples=101,
        declare function={normaldensity(\x)=exp(-\x^2/2)/sqrt(2*pi);},
        ]
        \addplot[draw=none,fill=black!12,domain=-4:-1.96] {normaldensity(x)} \closedcycle;
        \addplot[draw=none,fill=blue!35,domain=-1.96:1.96] {normaldensity(x)} \closedcycle;
        \addplot[draw=none,fill=black!12,domain=1.96:4] {normaldensity(x)} \closedcycle;
        \addplot[black,thick,smooth,domain=-4:4] {normaldensity(x)};
        \addplot[black,densely dashed] coordinates {(-1.96,0) (-1.96,{normaldensity(1.96)})};
        \addplot[black,densely dashed] coordinates {(1.96,0) (1.96,{normaldensity(1.96)})};
        \node[font=\small] at (axis cs:0,0.17) {$C=0.95$};
        \node[font=\small] (left tail) at (axis cs:-2.8,0.20) {$(1-C)/2=0.025$};
        \node[font=\small] (right tail) at (axis cs:2.8,0.20) {$(1-C)/2=0.025$};
        \draw[->] (left tail.south) -- (axis cs:-2.18,0.017);
        \draw[->] (right tail.south) -- (axis cs:2.18,0.017);
    \end{axis}
    \end{tikzpicture}
    \par\smallskip
    \small The cumulative area to the left of the axis position $z^*=1.960$ is the left tail plus the central region, so $\Phi(1.960)\approx 0.025+0.95=0.975$.
\end{minipage}

\clearpage

\subsection{Practice procedure}

\begin{enumerate}
    \item \textbf{Identify the estimates.} Name the population, the mean of interest, and its units. Use $\bar{x}$ as a single estimate of $\mu$ and the interval as a range of plausible values for $\mu$ under the model.
    \item \textbf{Calculate and compare intervals.} Multiply the table's $z^*$ by $\mathrm{SE}$, then use $\bar{x}\pm\mathrm{MOE}$ for the endpoints. Compare widths using $2\mathrm{MOE}$ before rounding the endpoints.
    \item \textbf{Measure observed coverage.} In a simulation with known $\mu$, count the intervals whose endpoints enclose $\mu$, including a match at either endpoint. Divide that count by the total.
    \item \textbf{Interpret confidence.} The nominal level is the level chosen for the procedure. It gives the approximate long-run coverage under the model. Observed coverage can vary across runs.
\end{enumerate}

\section{General Worked Problems}
\label{sec:general-worked}

\SectionProblemsTableOfContents{
    \SectionProblemLink{prob:c2-1}{Repeated estimates of market fees} \\
    \SectionProblemLink{prob:c2-2}{Two reports on household energy bills} \\
    \SectionProblemLink{prob:c2-3}{Bus commute times in a simulation}
}

\subsection{Repeated estimates of market fees}
\label{prob:c2-1}

\noindent
\begin{minipage}[t]{0.48\linewidth}
    \vspace{0pt}
    A city takes a random sample of market stalls to estimate the mean weekly stall fee. All fees are in thousand COP. Assume independent observations and normally distributed sample means with known $\mathrm{SE}=5/1.960$. The 95\% procedure uses $z^*=1.960$ and gives $\mathrm{MOE}=5$. The city's sample gives $\bar{x}=106$ and the interval $[101,111]$.

    The diagram shows the 95\% intervals from five fresh samples of the same size from a simulated market with mean $\mu=104$.
\end{minipage}
\hfill
\begin{minipage}[t]{0.48\linewidth}
    \vspace{0pt}
    \centering
    \begin{tikzpicture}
    \begin{axis}[
        width=\linewidth,height=4.6cm,axis lines=left,
        xmin=98,xmax=116,ymin=0.6,ymax=5.9,
        xlabel={Mean weekly fee (thousand COP)},ylabel={Simulated sample},
        ytick={1,2,3,4,5},xtick={99,104,109,114},
        tick label style={font=\small},label style={font=\small}]
        \addplot[black,dashed] coordinates {(104,0.6) (104,5.65)};
        \addplot[blue!65!black,thick,mark=|] coordinates {(101,1) (111,1)};
        \addplot[blue!65!black,thick,mark=|] coordinates {(99,2) (109,2)};
        \addplot[blue!65!black,thick,mark=|] coordinates {(103,3) (113,3)};
        \addplot[blue!65!black,thick,mark=|] coordinates {(105,4) (115,4)};
        \addplot[blue!65!black,thick,mark=|] coordinates {(102,5) (112,5)};
        \node[font=\small,anchor=east] at (axis cs:103.7,4.4) {$\mu=104$};
    \end{axis}
    \end{tikzpicture}
\end{minipage}

\QuestionsTasks[sec:general-worked]{
    {2}/{sol:c2-1-a}/{\textbf{Task 1.} \EstimationTask},
    {2}/{sol:c2-1-b}/{\textbf{Task 2.} Using the diagram, which option gives the correct number of intervals that cover the simulated mean and the correct interpretation of the city's 95\% confidence level?}
}

\begin{answeroptions}
    \item Four of five cover, so the observed 80\% is the city's long-run coverage rate.
    \item Four of five cover. The procedure covers the mean about 95\% of the time in the long run under the model. The city's calculated interval either covers its fixed mean or misses it.
    \item Four of five cover. The city's calculated interval has a 95\% probability of covering its fixed mean.
    \item Five of five cover. The procedure covers the mean about 95\% of the time in the long run under the model. The city's calculated interval either covers its fixed mean or misses it.
\end{answeroptions}

\clearpage
\noindent\begin{minipage}[t]{\linewidth}
\subsubsection*{Task 1 solution. Estimate the mean market fee}
\label{sol:c2-1-a}

\EstimationTable{
    Population parameter & $\mu$, the mean weekly stall fee for all market stalls in the city, in thousand COP. & The fixed, unknown mean the city wants to estimate. \\
    Point estimate & Reported sample mean $\bar{x}=106$ thousand COP. & A single estimate of the city's mean weekly stall fee. \\
    95\% interval & $\mathrm{MOE}=1.960(5/1.960)=5$ thousand COP. The endpoints are $[106-5,106+5]=[101,111]$. The width is $2(5)=10$ thousand COP. & A range of plausible values for the city's mean weekly stall fee under the model. \\
}

\secInicio[prob:c2-1][sec:general-worked]
\end{minipage}\par

\subsubsection*{Task 2 solution. Read the coverage diagram}
\label{sol:c2-1-b}

\textbf{Correct option B.} Each simulated interval is 10 thousand COP wide. Samples 1, 2, 3, and 5 cross the line at 104. Sample 4 starts at 105 and misses, so the observed coverage is $4/5=80\%$. The nominal level is 95\%. Over many fresh samples under the model assumptions, about 95\% of the resulting intervals would contain the fixed mean.

The claim of a 95\% probability for the city's interval fails. Its endpoints, 101 and 111, are fixed after sampling. The city's unknown mean either lies between them or falls outside them.

\secInicio[prob:c2-1][sec:general-worked]

\clearpage

\subsection{Two reports on household energy bills}
\label{prob:c2-2}

A utility takes a random sample of household accounts from one billing month. The sample mean is 120 thousand COP. Assume independent observations and a normal sampling model with known standard error 3.8 thousand COP. Both reports use the same sample and the rule $\bar{x}\pm\mathrm{MOE}$, with $\mathrm{MOE}=z^*\mathrm{SE}$.

\begin{center}
\begin{tabular}{@{}lcc@{}}
    \toprule
    Nominal confidence level & Critical value $z^*$ & Reported interval (thousand COP) \\
    \midrule
    90\% & 1.645 & $[114,126]$ \\
    99\% & 2.576 & $[110,130]$ \\
    \bottomrule
\end{tabular}
\end{center}

The reports round endpoints to the nearest thousand COP. The confidence levels refer to the procedures before rounding.

\QuestionsTasks[sec:general-worked]{
    {2}/{sol:c2-2-a}/{\textbf{Task 1.} \EstimationTask},
    {2}/{sol:c2-2-b}/{\textbf{Task 2.} Which option correctly compares the widths in the table and interprets the 99\% confidence level?}
}

\begin{answeroptions}
    \item The 99\% interval is wider, and it contains 99\% of individual bills.
    \item The 99\% interval is narrower, and the 99\% procedure covers the population mean about 99\% of the time in the long run under the model. The confidence level gives no percentage of individual bills within the reported endpoints.
    \item The 99\% interval is wider, and its reported endpoints have a 99\% probability of enclosing the population mean.
    \item The 99\% interval is wider, and the 99\% procedure covers the population mean about 99\% of the time in the long run under the model. The confidence level gives no percentage of individual bills within the reported endpoints.
\end{answeroptions}

\clearpage
\noindent\begin{minipage}[t]{\linewidth}
\subsubsection*{Task 1 solution. Interpret both reports}
\label{sol:c2-2-a}

\EstimationTable{
    Population parameter & $\mu$, the mean energy bill for all the utility's household accounts in that billing month, in thousand COP. & The fixed, unknown mean estimated in both reports. \\
    Point estimate & Reported sample mean $\bar{x}=120$ thousand COP. & A single estimate of the mean household bill for the month. \\
    90\% interval & $\mathrm{MOE}=1.645(3.8)=6.251$ thousand COP. The endpoints from $120\pm6.251$ are $113.749$ and $126.251$. The rounded interval is $[114,126]$ thousand COP. & The narrower reported range of plausible values for the mean household bill under the model. \\
    99\% interval & $\mathrm{MOE}=2.576(3.8)=9.7888$ thousand COP. The endpoints from $120\pm9.7888$ are $110.2112$ and $129.7888$. The rounded interval is $[110,130]$ thousand COP. & The wider reported range of plausible values for the same mean under the model. \\
}

\secInicio[prob:c2-2][sec:general-worked]
\end{minipage}\par

\subsubsection*{Task 2 solution. Compare confidence and width}
\label{sol:c2-2-b}

\textbf{Correct option D.} Before rounding, the widths are $2(6.251)=12.502$ and $2(9.7888)=19.5776$ thousand COP. The reported widths are $126-114=12$ and $130-110=20$ thousand COP. With the same standard error, the larger critical value for 99\% gives a larger $\mathrm{MOE}$ and puts the endpoints farther from 120.

Under the model assumptions, about 90\% and 99\% of the respective unrounded intervals would cover the fixed mean over repeated samples. Each calculated interval either covers the mean or misses it. The claim about 99\% of individual bills fails because coverage concerns the population mean. The proportion of bills from 110 to 130 thousand COP depends on the distribution of individual bills.

\secInicio[prob:c2-2][sec:general-worked]

\clearpage

\subsection{Bus commute times in a simulation}
\label{prob:c2-3}

A transit agency takes an independent random sample of 36 weekday passenger trips on one bus route. The sample mean is 34 minutes. Assume normally distributed trip times with known population standard deviation 12 minutes. For a 90\% interval the agency uses $z^*=1.645$ to find $\mathrm{MOE}$ and obtains $[30.71,37.29]$ minutes.

To study the 90\% procedure, the agency simulates 20 fresh samples of size 36 from a normal model with mean 35 minutes and the same standard deviation. The table groups the simulated intervals by sample number.

\begin{center}
\begin{tabular}{@{}lcccc@{}}
    \toprule
    Sample numbers & 1--5 & 6--10 & 11--15 & 16--20 \\
    \midrule
    Intervals containing 35 minutes & 5 & 4 & 3 & 5 \\
    \bottomrule
\end{tabular}
\end{center}

\QuestionsTasks[sec:general-worked]{
    {2}/{sol:c2-3-a}/{\textbf{Task 1.} \EstimationTask},
    {2}/{sol:c2-3-b}/{\textbf{Task 2.} Which option correctly totals the covering intervals in the table and assesses the claim that every run of 20 must have exactly 18 intervals that cover the mean?}
}

\begin{answeroptions}
    \item Seventeen of 20 cover. The claim fails because the 90\% procedure can give different counts in different runs of 20.
    \item Eighteen of 20 cover. The claim fails because the 90\% procedure can give different counts in different runs of 20.
    \item Seventeen of 20 cover. The claim fails because the observed 85\% is the procedure's long-run coverage rate.
    \item Seventeen of 20 cover. The claim fails because each calculated interval has a 90\% probability of containing the fixed mean.
\end{answeroptions}

\clearpage
\noindent\begin{minipage}[t]{\linewidth}
\subsubsection*{Task 1 solution. Estimate the mean trip time}
\label{sol:c2-3-a}

\EstimationTable{
    Population parameter & $\mu$, the mean duration of all weekday passenger trips on the route, in minutes. & The fixed, unknown mean the transit agency wants to estimate. \\
    Point estimate & Reported sample mean $\bar{x}=34$ minutes. & A single estimate of the mean weekday trip time on the route. \\
    90\% interval & $\mathrm{SE}=12/\sqrt{36}=2$ minutes and $\mathrm{MOE}=1.645(2)=3.29$ minutes. The endpoints of $34\pm3.29$ are $[30.71,37.29]$ minutes. The width is $2(3.29)=6.58$ minutes. & A range of plausible values for the route's mean weekday trip time under the model. \\
}

\secInicio[prob:c2-3][sec:general-worked]
\end{minipage}\par

\subsubsection*{Task 2 solution. Interpret a short run}
\label{sol:c2-3-b}

\textbf{Correct option A.} The table gives $5+4+3+5=17$ intervals that contain 35 minutes, for observed coverage of $17/20=85\%$. This run disproves the claim of exactly 18 covering intervals in every run.

The nominal level is 90\%. Under the model assumptions, about 90\% of intervals cover the fixed mean over many repetitions. This gives an average of about $0.90(20)=18$ covering intervals per run of 20, with counts that can vary from run to run. Each calculated interval either contains the mean or misses it.

\secInicio[prob:c2-3][sec:general-worked]

\clearpage

\section{Advanced Problems}
\label{sec:advanced}

\SectionProblemsTableOfContents{
    \SectionProblemLink{prob:c2-4}{Paired intervals for delivery times} \\
    \SectionProblemLink{prob:c2-5}{Sample size and produce spending}
}

\subsection{Paired intervals for delivery times}
\label{prob:c2-4}

A cooperative takes an independent random sample of 36 orders delivered in October. Their mean delivery time is 48 minutes. Assume normal delivery times with known population standard deviation 12 minutes. From that sample, the cooperative reports a 90\% interval of $[44.71,51.29]$ minutes and a 95\% interval of $[44.08,51.92]$ minutes. It uses $\bar{x}\pm\mathrm{MOE}$ with critical values $z^*=1.645$ and $z^*=1.960$, respectively.

In a simulation with a known mean and the same standard deviation, the cooperative computes a 90\% and a 95\% interval from each of 100 fresh samples of size 36. The row and column labels give the nominal levels. The cells count the observed outcomes for the paired intervals.

\begin{center}
\begin{tabular}{@{}lrr@{}}
    \toprule
    & 95\% interval covers & 95\% interval misses \\
    \midrule
    90\% interval covers & 89 & 0 \\
    90\% interval misses & 7 & 4 \\
    \bottomrule
\end{tabular}
\end{center}

\QuestionsTasks[sec:advanced]{
    {2}/{sol:c2-4-a}/{\textbf{Task 1.} \EstimationTask},
    {2}/{sol:c2-4-b}/{\textbf{Task 2.} Which option gives the observed coverage rates for the 90\% and 95\% procedures, in that order, and correctly uses the way each pair is constructed to judge whether a 90\% interval could cover the mean while its 95\% partner misses?}
}

\begin{answeroptions}
    \item 90\% and 95\%. The 95\% partner must cover whenever the 90\% interval covers because it is wider with the same centre.
    \item 89\% and 96\%. The 95\% partner can miss even when the 90\% interval covers because some 95\% intervals miss the mean.
    \item 89\% and 96\%. The 95\% partner must cover whenever the 90\% interval covers because it is wider with the same centre.
    \item 89\% and 89\%. The 95\% partner must cover whenever the 90\% interval covers because it is wider with the same centre.
\end{answeroptions}

\clearpage
\noindent\begin{minipage}[t]{\linewidth}
\subsubsection*{Task 1 solution. Interpret the two ranges}
\label{sol:c2-4-a}

\EstimationTable{
    Population parameter & $\mu$, the mean delivery time of all the cooperative's orders delivered in October, in minutes. & The fixed, unknown mean estimated by both intervals. \\
    Point estimate & Reported sample mean $\bar{x}=48$ minutes. & A single estimate of the mean October delivery time. \\
    90\% interval & $\mathrm{SE}=12/\sqrt{36}=2$ minutes and $\mathrm{MOE}=1.645(2)=3.29$ minutes. The endpoints of $48\pm3.29$ are $[44.71,51.29]$ minutes. The width is $2(3.29)=6.58$ minutes. & The narrower range of plausible values for the mean October delivery time under the model. \\
    95\% interval & With the same standard error, $\mathrm{MOE}=1.960(2)=3.92$ minutes. The endpoints of $48\pm3.92$ are $[44.08,51.92]$ minutes. The width is $2(3.92)=7.84$ minutes. & A wider range of plausible values for the same mean under the model, extending farther on both sides of 48. \\
}

\secInicio[prob:c2-4][sec:advanced]
\end{minipage}\par

\subsubsection*{Task 2 solution. Explain the coverage pattern}
\label{sol:c2-4-b}

\textbf{Correct option C.} The observed coverage for the 90\% procedure is $(89+0)/100=89\%$ from the covering row. For the 95\% procedure it is $(89+7)/100=96\%$ from the covering column.

Within each pair, the intervals share a centre and standard error. The larger critical value for 95\% gives a larger $\mathrm{MOE}$, so the 95\% interval contains the entire 90\% interval. The 95\% partner therefore covers whenever the 90\% interval covers, which accounts for the zero in the table.

Under the model assumptions, about 90\% and 95\% of the respective intervals would cover the fixed mean over repeated sampling. Each calculated interval either contains the mean or misses it.

\secInicio[prob:c2-4][sec:advanced]

\clearpage

\subsection{Sample size and produce spending}
\label{prob:c2-5}

A municipal survey estimates the mean amount households in one district spent on fresh produce during a specified week. Two teams take separate random samples. Assume independent household observations and normally distributed spending with known population standard deviation 50 thousand COP. Both teams use the 95\% rule $\bar{x}\pm\mathrm{MOE}$ with $z^*=1.960$.

\begin{center}
\begin{tabular}{@{}lrrc@{}}
    \toprule
    Team & Sample size & Sample mean (thousand COP) & 95\% interval (thousand COP) \\
    \midrule
    A & 25 & 205 & $[185.4,224.6]$ \\
    B & 100 & 205 & $[195.2,214.8]$ \\
    \bottomrule
\end{tabular}
\end{center}

The teams simulate 200 fresh samples at each size from a population with a known mean and the same standard deviation. Team A records 192 intervals that cover the mean. Team B records 187.

\QuestionsTasks[sec:advanced]{
    {2}/{sol:c2-5-a}/{\textbf{Task 1.} \EstimationTask},
    {2}/{sol:c2-5-b}/{\textbf{Task 2.} Which option correctly compares Team B's interval width with Team A's, gives the observed coverage rates for A and B, and assesses the claim that Team A's higher observed rate shows better long-run coverage?}
}

\begin{answeroptions}
    \item B's interval is half as wide. A covers 96\% and B covers 93.5\% in the simulations, so A has better long-run coverage.
    \item B's interval is half as wide. A covers 96\% and B covers 93.5\% in the simulations. Both procedures cover about 95\% of the time in the long run under the model.
    \item B's interval is one-quarter as wide. A covers 96\% and B covers 93.5\% in the simulations. Both procedures cover about 95\% of the time in the long run under the model.
    \item B's interval is half as wide. A and B each cover 95\% in the simulations and cover about 95\% of the time in the long run under the model.
\end{answeroptions}

\clearpage
\noindent\begin{minipage}[t]{\linewidth}
\subsubsection*{Task 1 solution. Estimate mean produce spending}
\label{sol:c2-5-a}

\EstimationTable{
    Population parameter & $\mu$, the mean spending on fresh produce among all households in the district during the specified week, in thousand COP. & The fixed, unknown district mean both teams want to estimate. \\
    Point estimate & Both samples report $\bar{x}=205$ thousand COP. & The same single estimate of the mean weekly produce spending from both samples. \\
    95\% interval (Team A) & $\mathrm{SE}=50/\sqrt{25}=10$ and $\mathrm{MOE}=1.960(10)=19.6$. The endpoints of $205\pm19.6$ are $[185.4,224.6]$ thousand COP. & The wider range of plausible values for the district mean under the model. \\
    95\% interval (Team B) & $\mathrm{SE}=50/\sqrt{100}=5$ and $\mathrm{MOE}=1.960(5)=9.8$. The endpoints of $205\pm9.8$ are $[195.2,214.8]$ thousand COP. & The narrower range of plausible values for the same district mean under the model. \\
}

\secInicio[prob:c2-5][sec:advanced]
\end{minipage}\par

\subsubsection*{Task 2 solution. Separate precision from coverage}
\label{sol:c2-5-b}

\textbf{Correct option B.} The widths are $2(19.6)=39.2$ and $2(9.8)=19.6$ thousand COP. Team B's fourfold sample size halves its standard error and therefore halves its $\mathrm{MOE}$ and interval width.

The observed rates are $192/200=96\%$ for A and $187/200=93.5\%$ for B. These rates can fluctuate across runs. Both procedures use the nominal level of 95\% and cover when the distance from $\bar{x}$ to $\mu$ is at most $\mathrm{MOE}=1.960\,\mathrm{SE}$. Under the model assumptions, that rule gives both procedures about 95\% coverage in the long run. Each calculated interval either contains the mean or misses it.

\secInicio[prob:c2-5][sec:advanced]

\end{document}
